%-------------------------
% Resume — Ivan Vanichkin
% Python Backend Developer
%-------------------------

\documentclass[a4paper,11pt]{article}

\usepackage[utf8]{inputenc}
\usepackage[T2A]{fontenc}
\usepackage[russian]{babel}
\usepackage{latexsym}
\usepackage[empty]{fullpage}
\usepackage{titlesec}
\usepackage[usenames,dvipsnames]{color}
\usepackage{enumitem}
\usepackage{fancyhdr}
\usepackage{tabularx}
\usepackage[hidelinks]{hyperref}

%----------PAGE STYLE----------

\pagestyle{fancy}
\fancyhf{}
\fancyfoot{}

\renewcommand{\headrulewidth}{0pt}
\renewcommand{\footrulewidth}{0pt}

%----------MARGINS----------

\addtolength{\oddsidemargin}{-0.375in}
\addtolength{\evensidemargin}{-0.375in}
\addtolength{\textwidth}{0.75in}
\addtolength{\topmargin}{-0.55in}
\addtolength{\textheight}{1.1in}

\urlstyle{same}
\raggedbottom
\raggedright
\setlength{\tabcolsep}{0in}

%----------SECTIONS----------

\titleformat{\section}
  {\vspace{-4pt}\scshape\raggedright\large}
  {}{0em}{}
  [\color{black}\titlerule \vspace{-5pt}]

%----------CUSTOM COMMANDS----------

\newcommand{\resumeItem}[2]{
  \item\small{
    \if\relax\detokenize{#1}\relax
      #2
    \else
      \textbf{#1}: #2
    \fi
    \vspace{-2pt}
  }
}

\newcommand{\resumeSubheading}[4]{
  \vspace{-1pt}\item
  \begin{tabular*}{0.97\textwidth}
    {l@{\extracolsep{\fill}}r}
    \textbf{#1} & #2 \\
    \textit{\small #3} & \textit{\small #4} \\
  \end{tabular*}
  \vspace{-5pt}
}

\newcommand{\resumeSubHeadingListStart}{
  \begin{itemize}[leftmargin=0pt,label={}]
}

\newcommand{\resumeSubHeadingListEnd}{
  \end{itemize}
}

\newcommand{\resumeItemListStart}{
  \begin{itemize}[leftmargin=18pt]
}

\newcommand{\resumeItemListEnd}{
  \end{itemize}\vspace{-5pt}
}

%===========================================
% DOCUMENT
%===========================================

\begin{document}

%===========================================
% HEADING
%===========================================

\noindent
\parbox[t]{\linewidth}{
    \vspace{3pt}

    {\Large \textbf{Иван Ваничкин}}
    \hfill
    \href{mailto:drdong654@gmail.com}
    {drdong654@gmail.com}
    \\[4pt]

    {\textbf{Junior Python Backend Developer}}
    \hfill
    +7~905~796~43~20
    \\[4pt]

    \href{https://github.com/drdong654}
    {github.com/drdong654}
    \hfill
    \href{https://t.me/VanichkinIM}
    {Telegram: @VanichkinIM}
}

\vspace{4pt}

%===========================================
% ABOUT
%===========================================

\section{О себе}

\small{
Junior Python Backend Developer с коммерческим опытом разработки
backend-сервисов и Telegram-приложений.
Переквалифицировался в Python backend-разработку через
профильные курсы, самостоятельную практику и коммерческую разработку.
Работаю с Python, FastAPI, PostgreSQL, SQLAlchemy,
asyncio и Docker. Проектирую бизнес-логику,
REST API и слой работы с данными, пишу автоматические тесты.
Имею практический опыт разработки RAG-систем
и интеграции LLM в backend-приложения.
}

%===========================================
% EXPERIENCE
%===========================================

\section{Опыт работы}

\resumeSubHeadingListStart

%----------VILDLY----------

\resumeSubheading
  {\Large Vildly}
  {\textbf{Python-разработчик}}
  {Backend-разработка}
  {Апрель 2026 --- настоящее время}

\vspace{7pt}

\textbf{Стек:}
Python, FastAPI, Aiogram 3, PostgreSQL, SQLAlchemy,
asyncpg, Docker, pytest, Git, RAG, pgvector

\vspace{3pt}

\resumeItemListStart

\resumeItem{}
{Разрабатываю backend-логику продукта на Python,
включая пользовательские сценарии, обработку данных
и взаимодействие компонентов приложения}

\resumeItem{}
{Разрабатываю REST API на FastAPI с разделением
API-слоя, бизнес-логики и слоя доступа к данным}

\resumeItem{}
{Работаю с PostgreSQL, SQLAlchemy и асинхронным
доступом к данным через asyncpg}

\resumeItem{}
{Реализовал регистрацию и авторизацию пользователей,
профили, валидацию входных данных и обработку ошибок}

\resumeItem{}
{Разрабатываю Telegram-интерфейс на Aiogram 3
и FSM-сценарии с интеграцией в backend приложения}

\resumeItem{}
{Работаю над RAG-системой и векторным поиском
по базе знаний с использованием PostgreSQL и pgvector}

\resumeItem{}
{Пишу unit- и integration-тесты на pytest,
контейнеризирую сервисы с Docker и работаю
с Git/GitHub в командном процессе разработки}

\resumeItemListEnd

\vspace{3pt}

\textbf{Результаты:}

\resumeItemListStart

\resumeItem{}
{Перевёл хранение пользовательских данных
с файлового прототипа на PostgreSQL}

\resumeItem{}
{Разделил Telegram-интерфейс, бизнес-логику
и слой доступа к данным}

\resumeItem{}
{Реализовал основные сценарии регистрации,
авторизации и управления пользовательским профилем}

\resumeItem{}
{Подготовил архитектурную основу для интеграции
RAG-поиска по базе знаний}

\resumeItemListEnd

%----------UPRAVA----------

\vspace{6pt}

\resumeSubheading
  {\Large Управа района Текстильщики}
  {\textbf{Ведущий специалист}}
  {Автоматизация внутренних процессов; совмещение}
  {Июль 2025 --- Июнь 2026}

\vspace{4pt}

\resumeItemListStart

\resumeItem{}
{Разрабатывал Telegram-ботов для автоматизации рекрутмента
и взаимодействия с участниками районных проектов}

\resumeItem{}
{Писал скрипты для обработки, структурирования и выгрузки
данных из БД, ускоряя анализ больших массивов рабочих данных}

\resumeItem{}
{Работал с корпоративными информационными системами
и внедрял цифровые инструменты во внутренние процессы}

\resumeItem{}
{Дорабатывал вёрстку сайта, обновлял контент
и контролировал актуальность публикуемой информации}

\resumeItemListEnd

\vspace{2pt}

\textbf{Результаты:}

\resumeItemListStart

\resumeItem{}
{Внедрил Telegram-ботов в реальные рабочие процессы
для автоматизации отдельных этапов рекрутмента}

\resumeItem{}
{Ускорил подготовку данных для анализа за счёт скриптов
обработки и выгрузки информации из БД}

\resumeItem{}
{Перевыполнил план по рекрутменту актива района --- более 100\%}

\resumeItemListEnd

\resumeSubHeadingListEnd

%===========================================
% HARD SKILLS
%===========================================

\section{Технические навыки}

\resumeItemListStart

\resumeItem{Backend}
{Python, FastAPI, Aiogram 3, asyncio, REST API, ООП}

\resumeItem{Databases}
{PostgreSQL, SQL, SQLAlchemy, asyncpg, pgvector}

\resumeItem{Testing}
{pytest, pytest-asyncio, unit testing,
integration testing, отладка}

\resumeItem{AI / RAG}
{RAG, LLM, Ollama, embeddings, vector search}

\resumeItem{Tools}
{Git, GitHub, Docker, Linux, Ruff, uv}

\resumeItemListEnd

%===========================================
% EDUCATION
%===========================================

\section{Образование}

\resumeSubHeadingListStart

\resumeSubheading
  {Реклама и связи с общественностью}
  {Российский новый университет}
  {Магистратура}
  {2025 --- 2027}

\vspace{2pt}

\resumeSubheading
  {Государственное и муниципальное управление}
  {Российский новый университет}
  {Бакалавриат}
  {2020 --- 2025}

\resumeSubHeadingListEnd

%===========================================
% ADDITIONAL EDUCATION
%===========================================

\section{Сертификаты и дополнительное обучение}

\resumeItemListStart

\resumeItem{Нетология, 2026}
{Курс «Основы Python: создаём телеграм-бота».
Сертификат № PYFREE-ASYNC-10158846; обучение в рамках
профессиональной переквалификации в Python backend-разработку}

\resumeItemListEnd

%===========================================
% LANGUAGES
%===========================================

\vspace{2pt}

\textbf{Языки:}
Русский --- родной, Английский --- B1

%===========================================

\end{document}
